\usepackage{xcolor}
\usepackage{listings}
\usepackage{inconsolata}
\usepackage{relsize}
\usepackage{amsmath,amssymb,array}
\usepackage{esz}
\usepackage{mathpartir}
\usepackage{graphicx}
\usepackage{xspace}
\usepackage{xurl}
\usepackage{caption}
\usepackage{bbding}
\usepackage{tabularx}

\hypersetup{colorlinks=true,linkcolor=blue,citecolor=blue,urlcolor=blue}
\def\UrlFont{\rmfamily}



\newcommand{\Section}[1]{\section{#1}}
\newcommand{\SubSection}[1]{\subsection{#1}}
\newcommand{\Paragraph}[1]{\vspace{0.5ex}\noindent\textbf{#1.}\xspace}

\newcommand{\TODO}[2]{%
    \noindent\textcolor{gray}{%
            {\scriptsize\textbf{TODO(#1)}: #2}
    }
}

\setlength{\belowcaptionskip}{-1ex}
\setlength{\textfloatsep}{0.2ex}
\newenvironment{Figure}{
    \begin{figure}[tb]
} {
    \end{figure}
}

% Full-width variant (spans both columns in the 2-column layout).
\newenvironment{WideFigure}{
    \begin{figure*}[!t]
} {
    \end{figure*}
}

% Full-width table variant (spans both columns in the 2-column layout).
\newenvironment{WideTable}{
    \begin{table*}[!t]
} {
    \end{table*}
}

% Abbreviations

\newcommand{\MVP}[0]{\textsf{MVP}\xspace}
\newcommand{\MSL}[0]{\textsf{MSL}\xspace}
\newcommand{\WP}[0]{\textsf{WP}\xspace}
\newcommand{\solidity}[0]{\textsf{Solidity}\xspace}
\newcommand{\kay}[0]{\textsf{K}\xspace}
\newcommand{\fstar}[0]{\textsf{F*}\xspace}
\newcommand{\verus}[0]{\textsf{Verus}\xspace}
\newcommand{\certora}[0]{\textsf{Certora}\xspace}


% IVL

% Based on esz.sty -- see CLAUDE.md

\newenvironment{ivl}{\zedindent=2ex\begin{zed}}{\end{zed}}

\newcommand{\requiresof}[2]{\Zpreop{requires\_of}\langle#1\rangle(#2)}
\newcommand{\abortsof}[2]{\Zpreop{aborts\_of}\langle#1\rangle(#2)}
\newcommand{\ensuresof}[2]{\Zpreop{ensures\_of}\langle#1\rangle(#2)}
\newcommand{\resultof}[2]{\Zpreop{result\_of}\langle#1\rangle(#2)}
\newcommand{\modifiesof}[1]{\Zpreop{modifies\_of}\langle#1\rangle}

% Variant-set names used in the implementation section.
\newcommand{\Clo}{\mathcal{C}}
\newcommand{\Par}{\mathcal{P}}
\newcommand{\Fld}{\mathcal{F}}

\newcommand{\Addr}{\Zpreop{address}}

\newcommand{\Mem}{\mathcal{M}}
\newcommand{\Mod}{\Delta}
\newcommand{\Type}{\mathcal{T}}
\newcommand{\Expr}{\mathcal{E}}

% Keywords for  IVL pseudocode
% ... esz already has \LET \IF \THEN \ELSE
\newcommand{\PROC}{\Zkeyword{proc}}
\newcommand{\FUN}{\Zkeyword{fun}}
\newcommand{\AXIOM}{\Zkeyword{axiom}}
\newcommand{\TRIGGER}{\Zinrel{trigger}}
\newcommand{\RETURNS}{\Zinrel{returns}}
\newcommand{\HAVOC}{\Zkeyword{havoc}}
\newcommand{\ASSUME}{\Zkeyword{assume}}
\newcommand{\ASSERT}{\Zkeyword{assert}}
\newcommand{\DATATYPE}{\Zkeyword{datatype}}
\newcommand{\IS}{\Zinrel{is}}
\newcommand{\TRUE}{\Zkeyword{true}}
\newcommand{\OLD}{\Zpreop{old}}
\newcommand{\BOOL}{\Zpreop{bool}}
\newcommand{\READ}{\Zpreop{read}}
\newcommand{\CALL}{\Zpreop{call}}
\newcommand{\VAR}{\Zkeyword{var}}
\newcommand{\FOREACH}{\Zkeyword{foreach}}

% State-labeled expression: \slabeled{S}{e} renders as S \sat e.
% \sat is a bar-then-tilde composite that reads like the literal |~.
% No standard single-glyph relation in amssymb / stmaryrd matches
% (closest is \vDash, which is |=). We therefore compose: \scriptstyle
% shrinks both glyphs; -1mu tightens the tilde against the bar while
% keeping a small visible gap.
% Kept in sync with the Move listings literate mapping below.
\newcommand{\sat}{\mathrel{{\scriptstyle\vert}\mkern-1mu{\scriptstyle\sim}}}
\newcommand{\slabeled}[2]{#1 \sat #2}

\newcommand{\sq}[1]{\overline{#1}}


% Source Transformation

\newcommand{\transform}[0]{\Large{$\leadsto$}}


%%%% Code

% Keep IEEEtran's default Times font; do not override with charter.
%\usepackage[charter]{mathdesign}
%\def\rmdefault{bch}
%\def\ttdefault{blg}

\lstdefinestyle{MoveStyle}{
    basicstyle=\ttfamily,
    keywordstyle=\color{black}, % don't bold keywords which is the default
    commentstyle=\color{gray}\normalfont\itshape,
    escapechar=@, % use to embed LaTeX into code
    breaklines=true,
    % \sat is set in \scriptstyle and followed by tight kerning, so its
    % right side-bearing is small; tail the replacement with \, so the
    % symbol doesn't abut the next token visually.
    literate={|~}{{$\sat$\,}}1
        {<=}{{$\leq$}}2
        {>=}{{$\geq$}}2
        {==}{{$=$}}2
        {!=}{{$\neq$}}2
        {==>}{{$\Longrightarrow$}}3
        {forall}{{$\forall$}}1
        {exists}{{$\exists$}}1,
}


\lstdefinelanguage{Move}{
    morekeywords={
        abort,
        aborts_if,
        aborts_of,
        acquires,
        address,
        as,
        assert,
        assume,
        borrow_global,
        borrow_global_mut,
        break,
        const,
        continue,
        copy,
        copyable,
        define,
        drop,
        else,
        ensures,
        ensures_of,
        exists,
        false,
        forall,
        friend,
        fun,
        global,
        has,
        havoc,
        if,
        in,
        include,
        invariant,
        key,
        let,
        loop,
        match,
        modifies,
        modifies_of,
        module,
        move,
        move_from,
        move_to,
        mut,
        native,
        num,
        old,
        onabort,
        pragma,
        proof,
        public,
        requires,
        requires_of,
        resource,
        result_of,
        return,
        schema,
        script,
        signer,
        spec,
        split,
        store,
        struct,
        true,
        u8,
        u64,
        u128,
        u256,
        update,
        use,
        with,
        where,
        while},
    sensitive=true,
    morecomment=[l]{//},
    morecomment=[s]{/*}{*/},
}

% Short inline delimiters for Move code: both |...| and !...! work.
% The ! alternative is useful when the snippet itself contains |,
% e.g. function-type syntax !|T1,...,Tn| -> T!.
\lstMakeShortInline[language=Move,style=MoveStyle]|
\lstMakeShortInline[language=Move,style=MoveStyle]!


% This defines a new environment for Move code.
\lstnewenvironment{Move}{
    \lstset{
        language=Move,
        style=MoveStyle,
        basicstyle=\relsize{-1}\ttfamily,
        keywordstyle=\color{blue},
    }
}{
}

% This defines a new environment for Move code in a box, without line numbers.
\lstnewenvironment{MoveBox}{
    \lstset{
        language=Move,
        style=MoveStyle,
        basicstyle=\relsize{-1}\ttfamily,
        keywordstyle=\color{blue},
    %numbers=left,
    %numberstyle=\scriptsize\color{gray},
    %frame=single,
    }
}{
}

% This defines a new environment for Move code in a box, with line numbers.
\lstnewenvironment{MoveBoxNumbered}{
    \lstset{
        language=Move,
        style=MoveStyle,
        basicstyle=\relsize{-1}\ttfamily,
        keywordstyle=\color{blue},
        numbers=left,
        numberstyle=\scriptsize\color{gray},
    %frame=single,
    }
}{
}

% This defines a new environment for diagnostics as produced by the prover.
\lstnewenvironment{MoveDiag}{
    \lstset{
        style=MoveStyle,
        basicstyle=\relsize{-2}\ttfamily,
    }
}{
}


%%% Local Variables:
%%% mode: latex
%%% TeX-master: "main"
%%% End:
